\documentclass[11pt,a4paper]{article}

\usepackage[T1]{fontenc}
\usepackage[utf8]{inputenc}
\usepackage[french]{babel}
\usepackage{amsmath,amssymb}
\usepackage[margin=2.3cm]{geometry}
\usepackage{booktabs}
\usepackage{array}
\usepackage{textcomp}
\usepackage{enumitem}
\usepackage{float}
\usepackage{longtable}
\usepackage{caption}
\usepackage{tikz}
\usetikzlibrary{arrows.meta,patterns}
\usepackage[colorlinks=true,linkcolor=blue,citecolor=blue,urlcolor=blue]{hyperref}

% Mêmes encadrés que le document de théorie du projet cartilage
\newsavebox{\implbox}
\newenvironment{implication}%
  {\par\medskip\begin{lrbox}{\implbox}\begin{minipage}{0.95\textwidth}%
   \textbf{Lien avec la pratique.}\itshape\ }%
  {\end{minipage}\end{lrbox}\noindent\fbox{\usebox{\implbox}}\par\medskip}

\newsavebox{\critbox}
\newenvironment{critere}%
  {\par\medskip\begin{lrbox}{\critbox}\begin{minipage}{0.95\textwidth}%
   \textbf{Critère de réussite.}\ }%
  {\end{minipage}\end{lrbox}\noindent\fbox{\usebox{\critbox}}\par\medskip}

\newcommand{\eps}{\varepsilon_r}
\newcommand{\ee}{\mathrm{e}}
\newcommand{\jj}{\mathrm{j}}

\title{\bfseries Tomographie micro-onde sur fantôme\\[2mm]
       \large Cadrage du banc : mesures, fantôme, reconstructions\\
       et spécification de la table XY}
\author{Notes de travail --- projet LiteVNA}
\date{\today}

\begin{document}
\maketitle

\begin{abstract}
\noindent
Ces notes fixent les quatre choix de départ du projet de tomographie, mené en
parallèle du projet cartilage avec le même LiteVNA ($1{,}4$--$6{,}3$~GHz). On
mesure à la fois en réflexion ($S_{11}$) et en transmission ($S_{21}$), avec deux
antennes face à face déplacées ensemble par une table XY pilotée par GRBL ;
l'objet reste fixe. Le premier fantôme est un cube de bois de $40$~mm de côté.
Les trois familles de reconstruction --- rétroprojection, méthode de Born,
inversion itérative --- sont développées sur une acquisition commune. La
spécification de la table en découle : ce n'est pas la finesse du pas qui la
dimensionne, mais la soustraction de la scène vide, qui exige une répétabilité
de position de l'ordre de $0{,}1$~mm.
\end{abstract}

\tableofcontents
\bigskip\hrule\bigskip

%=======================================================================
\section{Réflexion et transmission : les deux}
\label{sec:mesures}
%=======================================================================

Une image formée par une seule antenne déplacée devant l'objet, en réflexion ---
l'analogue du mode~B de l'échographie --- relève déjà de la tomographie en
réflexion. Mais la configuration la plus riche utilise deux antennes, et le
LiteVNA la mesure nativement :

\begin{itemize}
\item \textbf{l'antenne du port~1 émet et reçoit son écho} : c'est $S_{11}$, la
      réflexion ;
\item \textbf{l'antenne du port~2, en face, reçoit ce qui traverse} : c'est
      $S_{21}$, la transmission ;
\item les deux sont mesurés \textbf{dans le même balayage}. Le LiteVNA ne mesure
      ni $S_{22}$ ni $S_{12}$ : le port~2 ne fait que recevoir.
\end{itemize}

\paragraph{Pourquoi les deux se complètent.}
Dans l'approximation de Born, une mesure faite avec une onde incidente de
direction $\hat{\imath}$ et reçue dans la direction $\hat{s}$ renseigne la
composante de Fourier spatiale du contraste de permittivité au vecteur
\[
  \mathbf{K} = k\,(\hat{s} - \hat{\imath}), \qquad k = \frac{2\pi f}{c}.
\]
En réflexion ($\hat{s} = -\hat{\imath}$), $|\mathbf{K}| = 2k$ : ce sont les
hautes fréquences spatiales, c'est-à-dire les \textbf{interfaces} --- les faces
et les arêtes du cube. En transmission à faible angle
($\hat{s} \simeq \hat{\imath}$), $|\mathbf{K}|$ reste petit : ce sont les basses
fréquences spatiales, c'est-à-dire la \textbf{permittivité moyenne} traversée.
Une antenne seule ne remplit que la première partie de ce domaine.

C'est la configuration de l'holographie micro-onde en champ proche : deux
antennes face à face, balayées ensemble dans deux plans parallèles de part et
d'autre de l'objet, qui enregistrent réflexion et transmission sur une large
bande~\cite{amineh2011}.

\begin{figure}[H]
\centering
\begin{tikzpicture}[>=Stealth, font=\small]
  % axe des antennes
  \draw[dash dot, gray] (0,3.0) -- (0,-3.0);
  % étrier en C
  \draw[line width=2.2pt, gray!70] (-0.25,3.9) -- (-5.2,3.9) -- (-5.2,-3.9)
        -- (-0.25,-3.9);
  % antennes, ouverture tournée vers l'objet
  \filldraw[fill=blue!15, draw=blue!60!black]
        (-0.25,3.9) -- (0.25,3.9) -- (0.9,3.0) -- (-0.9,3.0) -- cycle;
  \filldraw[fill=blue!15, draw=blue!60!black]
        (-0.25,-3.9) -- (0.25,-3.9) -- (0.9,-3.0) -- (-0.9,-3.0) -- cycle;
  % plaque de polystyrène et son support fixe
  \fill[pattern=dots, pattern color=gray!70] (-4.4,-0.4) rectangle (4.6,-0.75);
  \draw[gray] (-4.4,-0.4) rectangle (4.6,-0.75);
  \draw[gray!70, line width=1.5pt] (4.6,-0.575) -- (5.4,-0.575) -- (5.4,-3.5);
  % cube
  \filldraw[fill=brown!35, draw=brown!60!black] (-0.4,-0.4) rectangle (0.4,0.4);
  % cotes
  \draw[gray, very thin] (0.9,3.0) -- (1.9,3.0) (0.4,0) -- (1.9,0)
        (0.9,-3.0) -- (1.9,-3.0);
  \draw[<->] (1.7,3.0) -- (1.7,0) node[midway, right] {$R \approx 150$~mm};
  \draw[<->] (1.7,0) -- (1.7,-3.0) node[pos=0.72, right] {$R$};
  % étiquettes
  \node[anchor=west] at (-0.1,4.35) {port 1 : émet et reçoit l'écho ($S_{11}$)};
  \node[anchor=west] at (-0.1,-4.35) {port 2 : reçoit la transmission ($S_{21}$)};
  \node[align=center, font=\footnotesize] at (-3.2,1.7)
        {étrier en C\\porté par la table XY};
  \node[anchor=east, font=\footnotesize] at (-0.5,0.25) {cube de bois, 40~mm};
  \node[font=\footnotesize] at (3.4,-1.05) {plaque de polystyrène};
  \node[anchor=west, font=\footnotesize] at (5.5,-2.6) {support fixe};
  % balayage
  \draw[<->, very thick, red!70!black] (-4.6,4.75) -- (-2.0,4.75)
        node[midway, above] {balayage X, et Y hors du plan};
\end{tikzpicture}
\caption{Configuration retenue, vue de côté. Les deux antennes, solidaires,
balayent ensemble ; le cube reste fixe sur une plaque de polystyrène
($\eps \simeq 1{,}02$) tenue hors de la zone balayée.}
\label{fig:config}
\end{figure}

\begin{implication}
\texttt{commun/instrument.py} ne rend aujourd'hui que $S_{11}$. Il faudra y
exposer $S_{21}$, que l'instrument fournit déjà, sans changer ce que voit le
projet cartilage.
\end{implication}

%=======================================================================
\section{Le fantôme : un cube en bois de 40 mm}
\label{sec:fantome}
%=======================================================================

Le premier fantôme est un cube de bois de $40$~mm de côté. \textbf{C'est un bon
choix, précisément parce que son contraste est faible} : la méthode de Born y
reste valable, donc les trois reconstructions du \S\ref{sec:reconstructions}
devraient donner le même résultat. Si elles concordent, la chaîne est validée.

Pour une épaisseur $L = 40$~mm et une longueur d'onde $\lambda$ dans l'air :
\[
  \Gamma = \frac{1-\sqrt{\eps}}{1+\sqrt{\eps}}, \qquad
  \Delta\tau = \frac{(\sqrt{\eps}-1)\,L}{c}, \qquad
  \Delta\varphi = 360^\circ\,\frac{(\sqrt{\eps}-1)\,L}{\lambda}.
\]

\begin{table}[H]
\centering
\begin{tabular}{lcccc}
\toprule
$\eps$ du bois & écho de la face & retard à travers 40~mm
  & phase à $6{,}3$~GHz & phase à $1{,}4$~GHz \\
\midrule
$1{,}5$ (très sec) & $-19{,}9$~dB & $30$~ps & $68^\circ$ & $15^\circ$ \\
$\mathbf{2{,}0}$ & $\mathbf{-15{,}3}$~\textbf{dB} & $\mathbf{55}$~\textbf{ps}
  & $\mathbf{125^\circ}$ & $\mathbf{28^\circ}$ \\
$3{,}0$ (humide) & $-11{,}4$~dB & $98$~ps & $222^\circ$ & $49^\circ$ \\
\bottomrule
\end{tabular}
\caption{Ce que voit le LiteVNA d'un cube de bois de $40$~mm.}
\label{tab:bois}
\end{table}

\paragraph{L'humidité compte.}
L'approximation de Born suppose un déphasage à travers l'objet inférieur à
$180^\circ$~\cite{kak1988}. À $\eps = 3$, on atteint $222^\circ$ en haut de bande
et Born décroche. Il faut donc un bois sec, droit de fil et sans nœud ; il est
aussi anisotrope, sa permittivité dépendant du sens du fil.

\paragraph{Le cube sera flou.}
$40$~mm valent $0{,}84\,\lambda$ à $6{,}3$~GHz. Le long de l'axe des antennes, la
résolution vaut $c/2B = 30{,}6$~mm, ou $52$~mm avec une fenêtre de Hann : autant
que le cube lui-même. Latéralement, elle vaut $15$ à $17$~mm
(\S\ref{sec:course}). On verra une tache de la bonne taille, pas des arêtes
nettes.

\paragraph{Il mesure lui-même sa permittivité.}
L'épaisseur étant connue, le retard en transmission donne directement
\[
  \eps = \left(1 + \frac{c\,\Delta\tau}{L}\right)^2, \qquad
  \frac{\mathrm{d}\eps}{\mathrm{d}\Delta\tau} = \frac{2\sqrt{\eps}\,c}{L},
\]
soit $0{,}02$ sur $\eps$ par picoseconde d'erreur, autour de $\eps = 2$.

\paragraph{La suite.}
Percer le cube et remplir le trou d'eau ($\eps \simeq 78$, environ $20$~mm de
diamètre) donnera un contraste fort. Born y échoue, et c'est là que l'inversion
itérative montrera ce qu'elle apporte.

%=======================================================================
\section{Les trois reconstructions}
\label{sec:reconstructions}
%=======================================================================

Les trois familles sont développées :

\begin{description}[style=nextline, leftmargin=1.2em]
\item[Rétroprojection (sommation cohérente)]
  Pour chaque point de l'image, on additionne les mesures corrigées du retard
  entre les antennes et ce point. Elle ne demande que les positions et la bande,
  aucun modèle du milieu : c'est la première à écrire, elle valide la géométrie.
\item[Méthode de Born (holographie, tomographie par diffraction)]
  Elle linéarise la relation entre le contraste de permittivité et le champ
  diffusé, et l'inverse dans le domaine de Fourier. Elle est quantitative tant
  que le contraste reste faible (\S\ref{sec:fantome}).
\item[Inversion itérative]
  Elle résout le problème non linéaire, par exemple par la méthode de Born
  itérative distordue ou par Gauss--Newton, en partant de la solution de Born.
  C'est la seule valable à fort contraste, la plus coûteuse, et la plus
  exigeante sur l'étalonnage.
\end{description}

Les trois peuvent partager \textbf{une seule acquisition}, à condition qu'elle
contienne, pour chaque position :

\begin{itemize}
\item $S_{11}$ et $S_{21}$ complexes sur toute la bande, enregistrés
      \textbf{bruts} ;
\item \textbf{le même balayage sans l'objet}, aux mêmes positions : Born et
      l'inversion itérative ont besoin de ce champ incident ;
\item les positions réellement atteintes.
\end{itemize}

%=======================================================================
\section{Spécification de la table XY}
\label{sec:table}
%=======================================================================

\subsection{Hypothèses}

\begin{itemize}
\item Les deux antennes sont face à face sur un \textbf{étrier rigide en C},
      porté par le chariot (figure~\ref{fig:config}).
\item Le cube est fixe entre elles, sur une plaque de polystyrène expansé.
\item Environ $R = 150$~mm entre chaque antenne et le centre du cube, au-delà du
      champ lointain d'une antenne de $50$~mm ($2D^2/\lambda = 105$~mm à
      $6{,}3$~GHz).
\item Bande $1{,}4$--$6{,}3$~GHz, donc $\lambda = 47{,}6$~mm en haut de bande.
\item Table XY à moteurs NEMA~17, commandée par GRBL.
\end{itemize}

\subsection{Ce qui dimensionne la table}
\label{sec:dimensionnement}

Une erreur de position $\delta$ le long de l'axe des antennes déphase la mesure
de
\[
  \delta\varphi_R = \frac{4\pi\,\delta}{\lambda} = 15{,}1^\circ/\text{mm}
  \quad\text{en réflexion}, \qquad
  \delta\varphi_T = \frac{2\pi\,\delta}{\lambda} = 7{,}6^\circ/\text{mm}
  \quad\text{en transmission}.
\]

Pour la \emph{focalisation} de l'image, c'est peu : une erreur de phase
quadratique moyenne de $22{,}5^\circ$ ne coûte que $0{,}67$~dB de gain, et
correspond à $1{,}5$~mm en réflexion. \textbf{Ce qui dimensionne la table, c'est
la soustraction de la scène vide.} Un écho parasite $B$ --- le bâti, le support
--- retranché avec une erreur de phase $\delta\varphi$ laisse un résidu
$|B|\,|1-\ee^{\jj\delta\varphi}| \simeq |B|\,\delta\varphi$. Pour garder ce résidu
sous $3\,\%$, la même cible que pour le projet cartilage, il faut
\[
  \delta \le \frac{0{,}03\,\lambda}{4\pi} = 0{,}11~\text{mm en réflexion},
  \qquad
  \delta \le \frac{0{,}03\,\lambda}{2\pi} = 0{,}23~\text{mm en transmission}.
\]

\begin{table}[H]
\centering
\begin{tabular}{cccc}
\toprule
erreur de & \multicolumn{2}{c}{erreur de phase} & résidu de soustraction \\
\cmidrule(lr){2-3}
position & réflexion & transmission & en réflexion \\
\midrule
$0{,}05$~mm & $0{,}76^\circ$ & $0{,}38^\circ$ & $1{,}3\,\%$ \\
$0{,}10$~mm & $1{,}51^\circ$ & $0{,}76^\circ$ & $2{,}6\,\%$ \\
$0{,}25$~mm & $3{,}78^\circ$ & $1{,}89^\circ$ & $6{,}6\,\%$ \\
$0{,}50$~mm & $7{,}57^\circ$ & $3{,}78^\circ$ & $13{,}2\,\%$ \\
$1{,}00$~mm & $15{,}1^\circ$ & $7{,}57^\circ$ & $26{,}4\,\%$ \\
$2{,}00$~mm & $30{,}3^\circ$ & $15{,}1^\circ$ & $52{,}8\,\%$ \\
\bottomrule
\end{tabular}
\caption{Effet d'une erreur de position à $6{,}3$~GHz.}
\label{tab:phase}
\end{table}

\subsection{Spécification}

\begin{longtable}{>{\raggedright\arraybackslash}p{3.3cm}
                  >{\raggedright\arraybackslash}p{5.0cm}
                  >{\raggedright\arraybackslash}p{6.3cm}}
\caption{Spécification de la table XY.}\label{tab:spec}\\
\toprule
Grandeur & Spécification & Pourquoi \\
\midrule
\endfirsthead
\toprule
Grandeur & Spécification & Pourquoi \\
\midrule
\endhead
\bottomrule
\endfoot
\textbf{Course utile \mbox{X × Y}} & $\ge 300 \times 300$~mm ; $400 \times 400$~mm si
  possible & résolution latérale à $6{,}3$~GHz : $16{,}8$~mm avec $300$~mm de
  course, $14{,}9$~mm avec $400$ (\S\ref{sec:course}) \\
\addlinespace
\textbf{Pas de balayage} & $10$~mm & $\le \lambda/4 = 11{,}9$~mm : pas de lobes de
  réseau, quel que soit l'angle \\
\addlinespace
\textbf{Répétabilité}, même sens d'approche & $\le \pm 0{,}05$~mm visé,
  $\pm 0{,}1$~mm au maximum & soustraction de la scène vide à $3\,\%$
  (\S\ref{sec:dimensionnement}) \\
\addlinespace
\textbf{Jeu à l'inversion} & $\le 0{,}05$~mm, \textbf{ou toujours approcher
  dans le même sens} & un balayage en serpentin inverse le sens une ligne sur
  deux \\
\addlinespace
\textbf{Exactitude sur la course} & $\le \pm 0{,}25$~mm, après étalonnage des
  pas/mm (\texttt{\$100}, \texttt{\$101}) & fidélité géométrique de l'image ;
  $3{,}8^\circ$ au pire, négligeable \\
\addlinespace
\textbf{Équerrage X/Y} & $\le 0{,}1^\circ$, soit $0{,}5$~mm sur $300$~mm &
  distorsion de l'image \\
\addlinespace
\textbf{Planéité} le long de l'axe des antennes & $\le \pm 0{,}2$~mm sur toute
  la course & répétable, donc retirée par la scène vide \\
\addlinespace
\textbf{Rigidité de l'étrier} : mouvement relatif des deux antennes &
  $\le 0{,}1$~mm, accélérations comprises & le champ incident en $S_{21}$ doit
  rester identique à chaque position ($0{,}76^\circ$) \\
\addlinespace
\textbf{Stabilisation} après l'arrêt & $< 0{,}05$~mm avant le balayage, en
  $\le 0{,}5$~s & $0{,}05$~mm de mouvement pendant les $0{,}25$~s d'un balayage
  valent $0{,}76^\circ$ \\
\addlinespace
\textbf{Résolution de commande} & $\le 0{,}05$~mm & facile : courroie GT2 à
  $20$~dents au $1/16$ de pas, $0{,}0125$~mm \\
\addlinespace
\textbf{Vitesse, accélération} & $\sim 3000$~mm/min (\texttt{\$110},
  \texttt{\$111}) ; $200$ à $500$~mm/s² (\texttt{\$120}, \texttt{\$121}) & pas
  critiques ; rester modéré pour ne pas faire osciller l'étrier \\
\addlinespace
\textbf{Prise d'origine} & fins de course, \texttt{\$22=1}, répétabilité
  $\le 0{,}05$~mm ; nouvelle prise d'origine en fin de campagne, écart
  $\le 0{,}1$~mm & GRBL fonctionne en boucle ouverte : c'est le seul moyen de
  détecter des pas perdus \\
\addlinespace
\textbf{Moteurs à l'arrêt} & \textbf{toujours alimentés} (\texttt{\$1=255}) &
  un moteur coupé retombe sur le pas entier le plus proche, jusqu'à
  $\pm 0{,}1$~mm sur GT2 ; si la vibration de maintien gêne, prendre des
  pilotes TMC2209 \\
\addlinespace
\textbf{Câbles coaxiaux} & phase stable à la flexion, $\le \pm 2^\circ$ à
  $6{,}3$~GHz sur le mouvement réel ; chaîne porte-câbles, rayon de courbure
  $\ge 10\times$ le diamètre & $2^\circ$ équivalent à $0{,}13$~mm en réflexion.
  \textbf{C'est le prix de déplacer les antennes plutôt que l'objet} : les
  câbles comptent autant que la mécanique \\
\addlinespace
\textbf{Environnement} & aucun métal du bâti dans les faisceaux ; support du
  cube en polystyrène ($\eps \simeq 1{,}02$) & échos parasites \\
\end{longtable}

\subsection{Course et résolution latérale}
\label{sec:course}

Avec une course $L$ vue à la distance $R$, l'ouverture angulaire totale $\theta$
et la résolution latérale en réflexion valent
\[
  \tan\frac{\theta}{2} = \frac{L}{2R}, \qquad
  \delta x = \frac{\lambda}{4\sin(\theta/2)} \;\ge\; \frac{\lambda}{4}
  = 11{,}9~\text{mm}.
\]

\begin{table}[H]
\centering
\begin{tabular}{cccc}
\toprule
course $L$ & demi-angle $\theta/2$ & $\delta x$ à $6{,}3$~GHz & \\
\midrule
$200$~mm & $33{,}7^\circ$ & $21{,}4$~mm & \\
$\mathbf{300}$~\textbf{mm} & $\mathbf{45{,}0^\circ}$ & $\mathbf{16{,}8}$~\textbf{mm}
  & minimum retenu \\
$400$~mm & $53{,}1^\circ$ & $14{,}9$~mm & si possible \\
$600$~mm & $63{,}4^\circ$ & $13{,}3$~mm & peu de gain \\
\bottomrule
\end{tabular}
\caption{Résolution latérale selon la course, pour $R = 150$~mm.}
\label{tab:course}
\end{table}

Au-delà de $400$~mm, le gain est faible, et l'antenne doit encore voir le cube à
$45$--$53^\circ$ hors de son axe : son diagramme de rayonnement borne l'ouverture
réellement utile. Le pas de $10$~mm reste sous $\lambda/4$ même dans ce cas ;
$\lambda/(4\sin 45^\circ) = 16{,}8$~mm suffirait si l'on se limitait à
$\pm 45^\circ$.

\subsection{Matériel : ordres de grandeur}

\begin{table}[H]
\centering
\begin{tabular}{lcc}
\toprule
transmission & répétabilité & jeu \\
\midrule
courroie GT2 & $0{,}05$ à $0{,}1$~mm & $0{,}1$ à $0{,}3$~mm selon la tension \\
vis trapézoïdale T8, écrou anti-jeu & $0{,}02$ à $0{,}05$~mm & --- \\
vis à billes SFU1204 & environ $0{,}01$~mm & --- \\
\bottomrule
\end{tabular}
\caption{Ordres de grandeur, à confirmer par le test T1.}
\label{tab:materiel}
\end{table}

La courroie suffit si l'on approche toujours chaque point dans le même sens.

\subsection{Déroulement d'un balayage}

\paragraph{Arrêt, mesure, marche.}
Jamais de mouvement pendant un balayage. Pour chaque position : envoyer
\texttt{G90 G0 X\ldots{} Y\ldots}, puis \texttt{G4 P0.5}. GRBL vide son tampon de
mouvements avant d'exécuter la temporisation, et ne répond \texttt{ok} au
\texttt{G4} qu'une fois le mouvement fini et la pause écoulée~\cite{grbl} : on
lance alors le balayage du LiteVNA. L'autre solution est d'interroger
\texttt{?} jusqu'à l'état \texttt{<Idle>}.

\paragraph{Durée.}
$31 \times 31 = 961$ positions à environ $2$~s chacune (déplacement, $0{,}5$~s
d'attente, quatre balayages moyennés), plus les retours de ligne : \textbf{environ
$35$~minutes}, et le double avec la scène vide.

\paragraph{Dérive.}
Sur une telle durée, il faut suivre la dérive du LiteVNA : revenir sur un point de
référence au début de chaque ligne et corriger les mesures de la ligne par ce
point.

\subsection{Tests de réception}

\begin{enumerate}[label=\textbf{T\arabic*}]
\item \textbf{Répétabilité et jeu.} Un comparateur au $1/100$~mm, $30$ approches
      du même point dans le même sens, puis dans les deux sens.
\item \textbf{Exactitude.} Un réglet ou un pied à coulisse sur toute la course :
      on règle les pas/mm, puis on vérifie l'exactitude.
\item \textbf{Stabilisation.} Le VNA comme capteur de vibration : une plaque
      métallique fixe, des balayages enchaînés juste après l'arrêt ; le temps que
      met la phase à se stabiliser donne l'attente nécessaire.
\item \textbf{Réception du banc.} Deux balayages complets de la scène vide,
      soustraits position par position.
\end{enumerate}

\begin{critere}
Au test T4, le résidu de soustraction reste sous $3\,\%$ à chaque position. Ce
seul test couvre à la fois la mécanique, les câbles et la dérive de
l'instrument.
\end{critere}

\subsection{Limite : un seul angle de vue}

Un balayage plan ne regarde le cube que dans une seule direction. Le long de l'axe
des antennes, la résolution ne vient que de la bande passante, soit $30$ à
$52$~mm, autant que le cube. Pour ajouter des angles de vue plus tard, la voie~Z
de GRBL peut piloter un plateau tournant, réglé en pas par degré, sans changer de
contrôleur.

%=======================================================================
\begin{thebibliography}{9}
%=======================================================================
\bibitem{amineh2011}
R.~K.~Amineh, M.~Ravan, A.~Khalatpour et N.~K.~Nikolova,
\og Three-dimensional near-field microwave holography using reflected and
transmitted signals \fg{}, \emph{IEEE Transactions on Antennas and
Propagation}, vol.~59, n\textsuperscript{o}~12, p.~4777--4789, 2011,
\url{https://doi.org/10.1109/TAP.2011.2165496}.

\bibitem{kak1988}
A.~C.~Kak et M.~Slaney, \emph{Principles of Computerized Tomographic Imaging},
IEEE Press, 1988, chapitre~6 : \og Tomographic imaging with diffracting
sources \fg{}.

\bibitem{grbl}
Grbl v1.1, documentation de l'interface,
\url{https://github.com/gnea/grbl/wiki/Grbl-v1.1-Interface}.
\end{thebibliography}

\end{document}
